\documentclass[9pt,a4paper]{article}
\usepackage[T1]{fontenc}
\usepackage{lmodern}
\usepackage{microtype}
\usepackage{amsmath,amssymb,mathtools}
\usepackage{geometry}
\usepackage{fancyhdr}
\usepackage{booktabs}
\usepackage{array}
\usepackage{float}
\usepackage{caption}
\usepackage[hidelinks]{hyperref}
\usepackage{xurl}
\usepackage{lastpage}
\geometry{top=1.35cm,bottom=1.40cm,left=1.85cm,right=1.85cm,headheight=13pt,headsep=0.22cm,footskip=0.65cm}
\setlength{\parindent}{0pt}
\setlength{\parskip}{0.35em}
\linespread{0.98}
\setlength{\emergencystretch}{2em}
\captionsetup{font=small,labelfont=bf,labelsep=period,justification=justified,singlelinecheck=false}
\hypersetup{hidelinks,pdftitle={Online Resource 4 - Supplementary Sections S4, S7, and S8},pdfauthor={Seyfullah Enes Kotil}}
\pagestyle{fancy}
\fancyhf{}
\fancyhead[L]{Online Resource 4}
\fancyhead[R]{Supplementary Sections S4, S7, and S8}
\renewcommand{\headrulewidth}{0.4pt}
\fancyfoot[C]{\thepage\ of \pageref{LastPage}}
\setlength{\tabcolsep}{4pt}
\renewcommand{\arraystretch}{1.0}

\newcommand{\ABM}{\mathrm{ABM}}
\newcommand{\HD}{\mathrm{HD}}

\begin{document}

\begin{center}
{\large\bfseries Online Resource 4}\\[0.3em]
{\bfseries Supplementary Sections S4, S7, and S8}\\[1.0em]
{\bfseries Overlapping roots: a transmission-tree framework for population-level contact-tracing dynamics}\\[0.8em]
Bulletin of Mathematical Biology\\[0.6em]
Seyfullah Enes Kotil\\
Boğaziçi University, Department of Molecular Biology and Genetics, Istanbul, Türkiye\\
Bahçeşehir University, School of Medicine, Department of Biophysics, Istanbul, Türkiye\\[0.4em]
Corresponding author: enesseyfullah.kotil@bogazici.edu.tr
\end{center}
\vspace{0.3em}\hrule\vspace{0.8em}

\textbf{Resource caption.} Analytical and numerical supplements for the implicit quasi-steady-state solution, the perturbative QSS construction, and the Section 6 technical results and large-population agent-based-model ensemble protocol.

\textbf{Contents.} S4 Analytical and numerical details for the QSS solution (Sect.~4). S7 Perturbative construction of the QSS approximation (Sect.~5). S8 Technical results and large-population ABM ensemble protocol for Section 6. References cited in this online resource.

Associated code and archives cited in Sections S4.3 and S7.5 are supplied in Online Resource 3 (\texttt{ESM\_3.zip}).

\section*{Supplementary Section S4: Analytical and numerical details for the QSS solution}

This supplementary section provides the admissibility argument, numerical evaluation procedure, and ABM-ensemble comparison protocol supporting Sect.~4 of the main text. The notation and main-text equation numbers are retained below.

\subsection*{S4.1 Existence and admissibility of the implicit QSS solution}

For $0<c\leq1$, define
\begin{equation}
F_{R,c}(u)=\sqrt{\frac{R}{c}}\;\frac{I_{R-cu+1}\!\left(2\sqrt{cR}\right)}{I_{R-cu}\!\left(2\sqrt{cR}\right)},
\qquad 0\leq u\leq1.
\tag{S4.1}
\end{equation}
The self-consistency condition in Eq.~(4.7) is $F_{R,c}(u)=u$. Since $R-cu>-1$ for $R>0$, $0<c\leq1$, and $0\leq u\leq1$, and since $2\sqrt{cR}>0$, the modified Bessel functions appearing in Eq.~(S4.1) are positive and continuous. Hence $F_{R,c}$ is continuous and positive on $[0,1]$.

For the positive decaying ratio sequence corresponding to a fixed trial value $u$, Eq.~(4.2) can be written as
\begin{equation}
r_k(u)=\frac{R}{R+k-cu+cr_{k+1}(u)}.
\tag{S4.2}
\end{equation}
At $k=1$, $F_{R,c}(u)=r_1(u)$. Therefore,
\begin{equation}
F_{R,c}(1)=\frac{R}{R+1-c+cr_2(1)}<1,
\tag{S4.3}
\end{equation}
where the inequality is strict because $r_2(1)>0$. Since $F_{R,c}(0)>0$, the continuous function $F_{R,c}(u)-u$ is positive at $u=0$ and negative at $u=1$. The intermediate value theorem therefore gives at least one self-consistent solution
\begin{equation}
m_1\in(0,1).
\tag{S4.4}
\end{equation}
At any such solution, the denominator in Eq.~(S4.2) exceeds $R$: indeed, $k-cm_1>0$ for every $k\geq1$, and $cr_{k+1}>0$. Hence
\begin{equation}
0<r_k<1,\qquad k\geq1.
\tag{S4.5}
\end{equation}
Since $m_k=m_{k-1}r_k$ and $m_0=1$, it follows that
\begin{equation}
1=m_0>m_1>m_2>\cdots>0.
\tag{S4.6}
\end{equation}
For $k\geq2$, $cm_1<1$ and $cr_{k+1}>0$, so Eq.~(S4.2) also gives
\begin{equation}
r_k<\frac{R}{R+k-1}.
\tag{S4.7}
\end{equation}
Consequently,
\begin{equation}
0<m_k=m_1\prod_{j=2}^{k}r_j<m_1\prod_{j=2}^{k}\frac{R}{R+j-1}\longrightarrow0\qquad(k\to\infty).
\tag{S4.8}
\end{equation}
Thus every positive self-consistent solution of Eq.~(4.7) defines a genealogically admissible descendant sequence.

When $c=0$, Eq.~(4.2) becomes explicit:
\begin{equation}
r_k(R,0)=\frac{R}{R+k},
\qquad
m_k(R,0)=\prod_{j=1}^{k}\frac{R}{R+j},
\tag{S4.9}
\end{equation}
and therefore $m_1(R,0)=R/(R+1)$. For each tested $c>0$, the scalar root was obtained independently. Across the tested parameter grid, the selected roots vary continuously from the no-tracing value and thereby identify the positive branch used in Sect.~4. This selection does not require a claim of global uniqueness for all mathematical roots of Eq.~(4.7).

\subsection*{S4.2 Numerical evaluation of the implicit solution}

The primary calculation solves the scalar residual
\begin{equation}
G_{R,c}(u)=F_{R,c}(u)-u
\tag{S4.10}
\end{equation}
for $u\in(0,1)$. At $c=0$, the explicit solution is $u=R/(R+1)$. For each $c>0$, the root was obtained independently by Brent's safeguarded method on $[0,1]$, using absolute and relative tolerances of $10^{-14}$. Across the tested grid, the resulting roots varied continuously from the no-tracing value and remained on the selected positive branch.

Modified-Bessel ratios were evaluated with exponentially scaled functions. Because the numerator and denominator in Eq.~(4.7) have the same argument, the common exponential factor cancels, avoiding overflow or underflow while preserving the required ratio (Amos 1974). For $c>0$, once $m_1$ was obtained, the descendant coordinates were calculated from
\begin{equation}
m_k=\left(\frac{R}{c}\right)^{k/2}\frac{I_{R-cm_1+k}\!\left(2\sqrt{cR}\right)}{I_{R-cm_1}\!\left(2\sqrt{cR}\right)},
\qquad k\geq1.
\tag{S4.11}
\end{equation}
At $c=0$, the explicit expression in Eq.~(4.9) was used.

As an independent consistency calculation, the continued fraction was evaluated by backward recursion. For a terminal depth $K$, set
\begin{equation}
r_{K+1}^{(K)}=0
\tag{S4.12}
\end{equation}
and compute
\begin{equation}
r_k^{(K)}=\frac{R}{R+k-cm_1+cr_{k+1}^{(K)}},
\qquad k=K,K-1,\ldots,1.
\tag{S4.13}
\end{equation}
The truncated self-consistency equation $r_1^{(K)}=m_1$ was solved separately at $K=40$ and $K=80$. Over the Figure~3 grid, the maximum fixed-point residual of the scaled-Bessel calculation was below $3\times10^{-15}$; the largest difference between the $K=40$ and $K=80$ values of $m_1$, $m_2$, and $m_3$ was below $10^{-15}$, and the largest difference between the scaled-Bessel and $K=80$ continued-fraction calculations was below $10^{-15}$. The verification script and cellwise numerical results accompany this section.

\subsection*{S4.3 Constant-pool ABM-ensemble comparison protocol}

The constant-pool simulations use
\begin{equation}
R\in\{1,1.5,2,4\},\qquad c\in\{0,0.25,0.5,0.75,1\},
\tag{S4.14}
\end{equation}
with dimensionless rates
\begin{equation}
\Gamma=1,\qquad \gamma=0,\qquad \gamma_c=1,\qquad p_f=c.
\tag{S4.15}
\end{equation}
The susceptible pool is held constant, so transmission occurs at total rate $RI$. Tracing is active from $\tau=0$. Each realization begins with 160 unrelated infectious roots, and 120 independent realizations are generated for every parameter pair over $0\leq\tau\leq3.25$.

For each active realization, the normalized descendant coordinate is
\begin{equation}
m_k(\tau)=\frac{M_k(\tau)}{I(\tau)},\qquad I(\tau)>0.
\tag{S4.16}
\end{equation}
A coordinate is left undefined after the realization no longer has an active infectious population; no stopped or terminal genealogy is carried forward.

The normalized coordinates were evaluated on the uniform output grid used to generate Fig.~3. Because the output spacing is uniform, equal weighting of the retained observations gives the corresponding discrete time average. For replicate $\ell$ and depth $k$, let $\mathcal T_{k,\ell}$ be the recorded times in the prespecified interval
\begin{equation}
3\leq\tau\leq3.25
\tag{S4.17}
\end{equation}
for which $I(\tau)>0$ and $m_k(\tau)$ is defined. The replicate-level summary is
\begin{equation}
m_{k,\ell}^{\ABM}=\frac{1}{|\mathcal T_{k,\ell}|}\sum_{\tau_j\in\mathcal T_{k,\ell}}m_k(\tau_j),
\tag{S4.18}
\end{equation}
provided $|\mathcal T_{k,\ell}|>0$. Thus, values after extinction are never imputed or carried forward; a replicate is excluded for depth $k$ only when it has no defined observation in the averaging interval.

Let $\mathcal V_k$ denote the replicates with a defined summary and $n_k=|\mathcal V_k|$. The signed difference plotted in Fig.~3a--c is
\begin{equation}
e^{\ABM}_{k,\mathrm{QSS}}=\frac{1}{n_k}\sum_{\ell\in\mathcal V_k}m_{k,\ell}^{\ABM}-m_k^{\mathrm{QSS}},
\qquad k=1,2,3.
\tag{S4.19}
\end{equation}
Percentile 95\% confidence intervals were obtained from 2,000 bootstrap resamples of complete replicate-level summaries; recorded time points within a realization were therefore never resampled as independent observations. For the trajectory panels, complete replicate trajectories were resampled and the ensemble mean was recalculated at each recorded time to obtain pointwise 95\% confidence intervals. The representative panels use $R=2$ and $c=1$; the pale-yellow region marks the averaging interval and the horizontal dashed line gives the selected QSS target.

\section*{Supplementary Section S7: Perturbative construction of the QSS approximation}

This supplementary section derives the perturbative construction summarized in Sect.~5 of the main text: the coefficient recursion, its local finite-depth basis, and the numerical evaluation of the resulting approximation errors. The notation and main-text equation numbers are retained below.

\subsection*{S7.1 Rescaled QSS recurrence}

At $c=0$, the QSS recurrence has the no-tracing solution
\begin{equation}
m_{k,0}=\prod_{j=1}^{k}\frac{R}{R+j},
\qquad k\geq1.
\end{equation}
Define
\begin{equation}
\widehat m_k=\frac{m_k}{m_{k,0}},
\qquad
b_k=\frac{R+1}{R+k},
\qquad
\chi_k=\frac{R+1}{R+k+1},
\qquad
\varepsilon=\frac{cR}{(R+1)^2}.
\tag{S7.1}
\end{equation}
Substitution into the QSS recurrence gives
\begin{equation}
\widehat m_k=\widehat m_{k-1}+\varepsilon b_k\left(\widehat m_1\widehat m_k-\chi_k\widehat m_{k+1}\right),
\qquad k\geq1.
\tag{S7.2}
\end{equation}
For $R>0$ and $0\leq c\leq1$,
\begin{equation}
0\leq\varepsilon\leq\frac{R}{(R+1)^2}\leq\frac14.
\tag{S7.3}
\end{equation}
The final inequality follows from $(R-1)^2\geq0$ and is sharp at $R=1$, $c=1$.

\subsection*{S7.2 Coefficient recursion and finite genealogical dependence}

Use the regular expansion
\begin{equation}
\widehat m_k=1+\sum_{n\geq1}a_{k,n}\varepsilon^n,
\qquad a_{k,0}=1,
\qquad a_{0,n}=0\;\;(n\geq1).
\tag{S7.4}
\end{equation}
Matching equal powers of $\varepsilon$ in Eq.~(S7.2) yields
\begin{equation}
a_{k,n}=a_{k-1,n}+b_k\left[\sum_{j=0}^{n-1}a_{1,j}a_{k,n-1-j}-\chi_ka_{k+1,n-1}\right],
\qquad n\geq1.
\tag{S7.5}
\end{equation}
The recursion couples adjacent descendant depths while lowering the perturbative order by one in the deeper term. Consequently,
\[
a_{1,p}\longrightarrow a_{2,p-1}\longrightarrow\cdots\longrightarrow a_{p,1}\longrightarrow a_{p+1,0}.
\]
The deepest required quantity is the known no-tracing coefficient $a_{p+1,0}=1$. Hence every fixed perturbative order is determined by a finite genealogical calculation.

For $k=1$,
\begin{align}
a_{1,1}&=\frac{1}{R+2},
\tag{S7.6}\\
a_{1,2}&=\frac{7+2R-R^2}{(R+2)^2(R+3)}.
\tag{S7.7}
\end{align}
Therefore,
\begin{equation}
m_1^{[p]}(R,c)=\frac{R}{R+1}\left[1+\sum_{n=1}^{p}a_{1,n}(R)\left(\frac{cR}{(R+1)^2}\right)^n\right].
\tag{S7.8}
\end{equation}
The formulas used in main-text Fig.~4 are
\begin{align}
m_1^{[0]}&=\frac{R}{R+1},
\tag{S7.9a}\\
m_1^{[1]}&=\frac{R}{R+1}\left[1+\frac{cR}{(R+1)^2(R+2)}\right],
\tag{S7.9b}\\
m_1^{[2]}&=\frac{R}{R+1}\left[1+\frac{cR}{(R+1)^2(R+2)}+\frac{7+2R-R^2}{(R+2)^2(R+3)}\left(\frac{cR}{(R+1)^2}\right)^2\right].
\tag{S7.9c}
\end{align}

\subsection*{S7.3 Local finite-depth basis}

Consider a finite truncation retaining $\widehat m_1,\ldots,\widehat m_K$ and an analytic terminal prescription $\widehat m_{K+1}=T_K(\varepsilon)$ satisfying $T_K(0)=1$. The truncated QSS equations are analytic in the retained moments and $\varepsilon$. At $\varepsilon=0$, Eq.~(S7.2) reduces to $\widehat m_k=\widehat m_{k-1}$, so the no-tracing solution is $\widehat m_k=1$. The Jacobian with respect to the retained moments is lower bidiagonal with unit diagonal and is therefore nonsingular. Standard analytic implicit-function and regular-perturbation arguments consequently give a locally unique analytic branch whose coefficients satisfy Eq.~(S7.5) (Krantz and Parks 2002; Murdock 1999).

This is a local statement for each fixed finite truncation. The main-text claim does not require a global convergence theorem for the complete infinite series; instead, the accuracy of the explicit low-order formulas is assessed directly against the selected QSS branch.

\subsection*{S7.4 Numerical evaluation for main-text Figure 4}

The selected QSS branch was evaluated by solving the scalar fixed-point equation generated by a backward continued fraction. Main-text Fig.~4 was generated using continued-fraction depth 80 and a bracketed scalar solve on $0\leq m_1\leq1$. The error curves were evaluated at $R\in\{1,1.5,2,4\}$ and at 1001 equally spaced values of $c$ on $[0,1]$. For retained order $p$,
\begin{equation}
e_p^{\mathrm{pert}}(R,c)=\left|m_1^{[p]}(R,c)-m_1^{\mathrm{QSS}}(R,c)\right|.
\end{equation}
The maximum errors over the $c$-grid are summarized in Table~S7.1.

\begin{table}[H]
\centering
\caption{Maximum absolute perturbative errors $e_p^{\mathrm{pert}}$ over $c\in[0,1]$ for the parameter values used in main-text Fig.~4.}
\label{tab:S7-1}
\small
\begin{tabular}{@{}rccc@{}}
\toprule
$R$ & $p=0$ & $p=1$ & $p=2$ \\
\midrule
1.0 & $5.0374326\times10^{-2}$ & $8.7076597\times10^{-3}$ & $1.7632153\times10^{-3}$ \\
1.5 & $4.6553438\times10^{-2}$ & $5.4105804\times10^{-3}$ & $5.5180492\times10^{-4}$ \\
2.0 & $3.9907541\times10^{-2}$ & $2.8705039\times10^{-3}$ & $1.0154532\times10^{-5}$ \\
4.0 & $2.1111821\times10^{-2}$ & $2.2151187\times10^{-4}$ & $1.4024202\times10^{-4}$ \\
\bottomrule
\end{tabular}
\end{table}

For every displayed $R$ and $p$, the maximum occurred at $c=1$. The signed $p=2$ error for $R=2$ crossed zero at $c=0.8942672386$.

The QSS solver was checked over $R\in[0.05,10]$ and $c\in[0.01,1]$ on a $124\times100$ diagnostic grid. Depths 40, 80, and 160 agreed to double-precision resolution. At depth 160, the maximum fixed-point residual was $2.33\times10^{-15}$, and the maximum discrepancy between the continued-fraction evaluation and the evaluation based on the modified-Bessel representation was $4.44\times10^{-15}$.

\subsection*{S7.5 Reproducibility files}

The accompanying archive provides the complete main-text Fig.~4 curve data, the maximum-error summary, the symbolic coefficient audit, the QSS solver diagnostics, and standalone generation scripts. These materials are stored under \texttt{data/}, \texttt{diagnostics/}, and \texttt{reproducibility/}; the symbolic audit independently verifies the displayed coefficients.

\section*{S8 Technical results and large-population ABM ensemble protocol for Section 6}

\subsection*{S8.1 Equality of the growth and replacement thresholds}

On the selected constant-pool QSS branch, the normalized moments satisfy
\begin{equation}
0=R(m_{k-1}-m_k)-km_k+c(m_1m_k-m_{k+1}),
\qquad k\geq1,
\label{eq:S8-qss-recurrence}
\end{equation}
with $m_0=1$. Define
\begin{equation}
B(R,c)=1+\sum_{k\geq1}\frac{(-c)^k}{(k+1)!}\,m_k(R,c).
\label{eq:S8-B-definition}
\end{equation}
The factorial weights make all sums below absolutely convergent. Multiply Eq.~\eqref{eq:S8-qss-recurrence} by $(-c)^k/(k+1)!$ and sum over $k\geq1$. After index shifts, the transmission terms generate the series in $\mathcal R_{\mathrm{time}}$, while the remaining terms collect into $gB$, giving
\begin{equation}
\mathcal R_{\mathrm{time}}(R,c)-1
=g(R,c)B(R,c).
\label{eq:S8-threshold-identity}
\end{equation}
For completeness, the weighted sum can be written directly as
\begin{align}
0
&=\sum_{k\geq1}\frac{(-c)^k}{(k+1)!}
\left[R(m_{k-1}-m_k)-km_k+c(m_1m_k-m_{k+1})\right] \\
&=\mathcal R_{\mathrm{time}}-1-gB.
\label{eq:S8-weighted-sum}
\end{align}
Moreover, Eq.~\eqref{eq:S8-B-definition} is an alternating series with strictly decreasing term magnitudes because
\begin{equation}
\frac{c^{k+1}m_{k+1}/(k+2)!}{c^km_k/(k+1)!}
=\frac{c}{k+2}\frac{m_{k+1}}{m_k}
\leq\frac{1}{k+2}<1.
\end{equation}
Hence
\begin{equation}
B(R,c)\geq1-\frac{cm_1}{2}>\frac12.
\label{eq:S8-B-positive}
\end{equation}
Therefore
\begin{equation}
g(R,c)=0
\quad\Longleftrightarrow\quad
\mathcal R_{\mathrm{time}}(R,c)=1,
\label{eq:S8-threshold-equivalence}
\end{equation}
although the two quantities have different numerical values and interpretations away from the boundary.

\subsection*{S8.2 Finite-depth evaluation of $\mathcal R_{\mathrm{time}}$}

Retaining genealogical coordinates through $m_{K-1}$, define
\begin{equation}
\mathcal R_{\mathrm{time}}^{(K)}(R,c)
=R\left[
1+\sum_{k=0}^{K-1}\frac{(-c)^{k+1}}{(k+2)!}\,m_k(R,c)
\right].
\label{eq:S8-Rtime-K}
\end{equation}
Because $m_k$ is nonincreasing and $0\leq c\leq1$, the omitted tail is alternating with decreasing magnitude. Its error is bounded by the first omitted term:
\begin{equation}
\left|
\mathcal R_{\mathrm{time}}-
\mathcal R_{\mathrm{time}}^{(K)}
\right|
\leq
R\frac{c^{K+1}}{(K+2)!}\,m_K
\leq Rm_K.
\label{eq:S8-Rtime-tail}
\end{equation}
The factorial factor makes the lifetime replacement quantity substantially less sensitive to deep-tail error than an unweighted moment sum.

\subsection*{S8.3 Propagation of QSS approximation errors}

If an approximation $m_1^{[p]}$ is used in the growth rate, then
\begin{equation}
g^{[p]}-g^{\mathrm{QSS}}
=-c\left(m_1^{[p]}-m_1^{\mathrm{QSS}}\right).
\label{eq:S8-growth-error}
\end{equation}
If approximate moments $m_k^{[p]}$ are used through depth $K-1$, then
\begin{align}
\left|
\mathcal R_{\mathrm{time}}^{[p,K]}
-
\mathcal R_{\mathrm{time}}^{\mathrm{QSS}}
\right|
&\leq
R\sum_{k=1}^{K-1}
\frac{c^{k+1}}{(k+2)!}
\left|m_k^{[p]}-m_k^{\mathrm{QSS}}\right| \\
&\quad+
R\frac{c^{K+1}}{(K+2)!}m_K^{\mathrm{QSS}}.
\label{eq:S8-Rtime-error}
\end{align}
The first term is the error in the retained coordinates and the second is the unresolved genealogical tail. This analysis concerns approximation of the selected QSS branch; it is distinct from finite-population ABM--QSS discrepancy.

\subsection*{S8.4 Large-population constant-pool ABM ensemble protocol}

The ABM uses the same transmission, primary-removal, detection, and instantaneous one-step forward-tracing events as the deterministic hierarchy. The susceptible pool is held constant in the transmission rate. The decomposition used for main-text Fig.~5 is
\begin{equation}
\frac{\gamma}{\Gamma}=0,
\qquad
\frac{\gamma_c}{\Gamma}=1,
\qquad
p_f=c.
\label{eq:S8-abm-decomposition}
\end{equation}
Thus, the deterministic prediction depends on $R$ and $c$, while the ABM realizes the corresponding finite-population event system.

\begin{table}[H]
\centering
\caption{Simulation design used for main-text Fig.~5.}
\label{tab:S8-simulation-design}
\scriptsize
\begin{tabular}{p{0.18\textwidth}p{0.29\textwidth}p{0.43\textwidth}}
\toprule
Component & Cells & Design \\
\midrule
Panels (a)--(b) & $R=1,1.5,2,4$; $c=0,0.25,0.5,0.75,1$ & 300 independent paths per cell; prespecified burn-in and windows; whole-replicate bootstrap intervals. \\
Cohort follow-up & Same 20 cells & Complete follow-up of every entrant; no lifetime or required production interval is censored. \\
Panel (c) & $c=0.25,0.5,0.75,1$; two prespecified $R$-bracketing cells & Pooled growth estimates, linear interpolation, and pointwise whole-replicate bootstrap intervals. \\
\bottomrule
\end{tabular}
\end{table}

For replicate $r$, let
\begin{equation}
P_r=\int_{\tau_{0,r}}^{\tau_{1,r}}I_r(\tau)\,d\tau
\label{eq:S8-parent-person-time}
\end{equation}
be the parent infectious person-time in the prespecified measurement or cohort-entry window. If $N_{+,r}$ is the number of new infections and $N_{-,r}$ is the number of infectious individuals removed by primary removal or tracing during the growth-measurement window, the pooled event/person-time estimator is
\begin{equation}
\widehat g
=
\frac{\sum_r(N_{+,r}-N_{-,r})}{\sum_rP_r}.
\label{eq:S8-g-estimator}
\end{equation}
Here $N_{-,r}$ counts removed individuals rather than removal-event occurrences, so tracing events that remove more than one infectious individual contribute their full population change.

Let $\mathcal C_r$ be the set of infections entering during the cohort window, and let $L_{jr}$ be the complete remaining infectious lifetime of cohort member $j$. The pooled infectious-time replacement estimator is
\begin{equation}
\widehat{\mathcal R}_{\mathrm{time}}
=
\frac{\displaystyle\sum_r\sum_{j\in\mathcal C_r}L_{jr}}
{\displaystyle\sum_rP_r}.
\label{eq:S8-Rtime-estimator}
\end{equation}
The event system is continued along each path until every enrolled cohort member has been removed. The entry window, burn-in rule, and production sample size are fixed before the confirmatory runs.

Confidence intervals are obtained by resampling complete replicates and recomputing the pooled ratio in Eqs.~\eqref{eq:S8-g-estimator} and~\eqref{eq:S8-Rtime-estimator}. Serial observations within a replicate are not treated as independent. For panel (c), let $(R^-,\widehat g^-)$ and $(R^+,\widehat g^+)$ denote the two prespecified cells bracketing zero. The interpolated estimate is
\begin{equation}
\widehat R_*(c)
=R^-+
\frac{-\widehat g^-}{\widehat g^+-\widehat g^-}
(R^+-R^-),
\label{eq:S8-threshold-estimator}
\end{equation}
and each bootstrap resample recomputes the two pooled growth estimates and the interpolation.

\subsection*{S8.5 Analytical and numerical checks}

The analytical curves use the same selected-QSS solver as main-text Figs.~4 and 5. The final audit gave a maximum fixed-point residual of $1.887\times10^{-13}$, a maximum accepted depth-doubling difference of $1.943\times10^{-13}$, and a maximum modified-Bessel/continued-fraction difference of $1.894\times10^{-13}$. The control-boundary values used in Fig.~5c are
\begin{equation}
\begin{array}{c|ccccc}
c & 0 & 0.25 & 0.50 & 0.75 & 1.00\\ \hline
R_*(c) & 1.00000000 & 1.13568648 & 1.29328143 & 1.47215210 & 1.66978036
\end{array}.
\label{eq:S8-critical-values}
\end{equation}
The threshold identity was checked on 8,888 parameter points. Across the 20 cohort cells, the maximum relative Monte Carlo standard error was 0.50\%, the maximum analytical discrepancy was 0.67\%, and 19 of 20 pointwise intervals contained their targets; all panel-(c) intervals contained the analytical critical curve. These comparisons assess large-population ensemble behavior of the same event system, not an unrelated stochastic model.

\section*{References}

The continued-fraction and special-function numerical procedures in Section~S4.2 use Gautschi (1967) and Olver et al. (2010).

\begin{list}{}{\setlength{\leftmargin}{1.2em}\setlength{\itemindent}{-1.2em}\setlength{\itemsep}{0.4em}}
\item Amos DE (1974) Computation of modified Bessel functions and their ratios. Math Comput 28:239--251. \url{https://doi.org/10.1090/S0025-5718-1974-0333287-7}
\item Gautschi W (1967) Computational aspects of three-term recurrence relations. SIAM Rev 9:24--82. \url{https://doi.org/10.1137/1009002}
\item Krantz SG, Parks HR (2002) The implicit function theorem: history, theory, and applications. Birkhäuser, Boston. \url{https://doi.org/10.1007/978-1-4612-0059-8}
\item Murdock JA (1999) Perturbations: theory and methods. Classics in Applied Mathematics, vol 27. SIAM, Philadelphia. \url{https://doi.org/10.1137/1.9781611971095}
\item Olver FWJ, Lozier DW, Boisvert RF, Clark CW (eds) (2010) NIST handbook of mathematical functions. Cambridge University Press, New York. \url{https://www.nist.gov/publications/nist-handbook-mathematical-functions}
\end{list}

\end{document}
